The source code is organized in two main directories, separating the frontend logic from the backend logic.

\subsection{Backend}
\begin{verbatim}
/
├── src
│   ├── main
│   │   ├── java/it/polimi/se/bbp
│   │   │   ├── config
│   │   │   ├── controller
│   │   │   ├── dto
│   │   │   │   ├── mapbox
│   │   │   │   ├── openmeteo
│   │   │   │   ├── request
│   │   │   │   ├── response
│   │   │   │   └── result
│   │   │   ├── entity
│   │   │   ├── enums
│   │   │   ├── exception
│   │   │   ├── geo
│   │   │   ├── mapper
│   │   │   ├── repository
│   │   │   ├── security
│   │   │   ├── service
│   │   │   └── specification
│   │   └── resources
│   │       └── application.properties
│   └── test
└── pom.xml
\end{verbatim}

\begin{description}
    \item[config] Contains Spring configuration classes, and external services such as Mapbox and OpenMeteo.
    \item[controller] Contains REST Controllers that expose API endpoints, handle HTTP requests, and delegate logic to services.
    \item[dto] Collects Data Transfer Objects, separated into request and response, as well as DTOs specific to external APIs.
    \item[entity] Defines JPA entities that map database tables.
    \item[exception] Centralized exception handling and custom exception definitions.
    \item[geo] Utility classes for geospatial calculations.
    \item[mapper] Interfaces for automatic conversion between Entities and DTOs.
    \item[repository] Interfaces for database interaction.
    \item[security] Components for JWT authentication.
    \item[service] Core business logic of the application.
    \item[specification] Classes for constructing dynamic queries used for search filters.
\end{description}

\subsection{Frontend}
\begin{verbatim}
/
├── public
├── src
│   ├── api
│   ├── assets
│   ├── components
│   ├── composables
│   ├── config
│   ├── constants
│   ├── router
│   ├── services
│   ├── stores
│   ├── styles
│   ├── types
│   ├── utils
│   ├── views
│   ├── App.vue
│   └── main.ts
├── index.html
├── package.json
├── tailwind.config.ts
└── vite.config.ts
\end{verbatim}

\begin{description}
    \item[api] HTTP client configuration and request interceptor management.
    \item[components] Reusable Vue components in the interface.
    \item[composables] Functions encapsulating reusable stateful logic.
    \item[router] Client-side routing configuration defining page navigation.
    \item[services] Modules interfacing the frontend with backend APIs.
    \item[stores] Global application state management via Pinia.
    \item[styles] Global CSS files.
    \item[types] Shared TypeScript interface definitions to ensure strong data typing.
    \item[views] Main application pages that assemble the various components.
    \item[utils] Pure utility functions for validation, date formatting, and static calculations.
\end{description}